\clearpage\phantomsection\label{exercise-index-4}
\BookHeading{Percorso degli esercizi guidati}
Quindici esercizi accompagnano lo studio del modulo. Ogni traccia segue le sezioni richieste e rinvia a uno svolgimento completo. Si suggerisce di formulare la previsione e tentare i passaggi prima di consultare la soluzione.
\begin{longtable}{@{}p{13mm}p{75mm}r@{}}\toprule Codice & Argomento & Pagina\\\midrule\endhead
\hyperref[ex-G4.01]{G4.01} & Ricostruire l'interferenza in Deutsch--Jozsa e Bernstein--Vazirani & \pageref{ex-G4.01}\\[3pt]
\hyperref[ex-G4.02]{G4.02} & Recuperare il periodo binario di Simon & \pageref{ex-G4.02}\\[3pt]
\hyperref[ex-G4.03]{G4.03} & Stimare una sovrapposizione mediante SWAP test & \pageref{ex-G4.03}\\[3pt]
\hyperref[ex-G4.04]{G4.04} & Ricostruire Grover come due riflessioni & \pageref{ex-G4.04}\\[3pt]
\hyperref[ex-G4.05]{G4.05} & Seguire una ricerca con budget e soglia variabile & \pageref{ex-G4.05}\\[3pt]
\hyperref[ex-G4.06]{G4.06} & Costruire una Fourier a quattro livelli & \pageref{ex-G4.06}\\[3pt]
\hyperref[ex-G4.07]{G4.07} & Leggere un'autofase esatta & \pageref{ex-G4.07}\\[3pt]
\hyperref[ex-G4.08]{G4.08} & Calcolare una fase non rappresentabile esattamente & \pageref{ex-G4.08}\\[3pt]
\hyperref[ex-G4.09]{G4.09} & Dalle autofasi di Grover al conteggio & \pageref{ex-G4.09}\\[3pt]
\hyperref[ex-G4.10]{G4.10} & Costruire RSA e recuperare la chiave dai fattori & \pageref{ex-G4.10}\\[3pt]
\hyperref[ex-G4.11]{G4.11} & Verificare somma e moltiplicazione modulari reversibili & \pageref{ex-G4.11}\\[3pt]
\hyperref[ex-G4.12]{G4.12} & Costruire gli autostati di un'orbita modulare & \pageref{ex-G4.12}\\[3pt]
\hyperref[ex-G4.13]{G4.13} & Ricostruire e verificare un denominatore & \pageref{ex-G4.13}\\[3pt]
\hyperref[ex-G4.14]{G4.14} & Svolgere Shor per \(N\) \(=\) 15 & \pageref{ex-G4.14}\\[3pt]
\hyperref[ex-G4.15]{G4.15} & Svolgere Shor per \(N\) \(=\) 21 & \pageref{ex-G4.15}\\[3pt]
\bottomrule\end{longtable}
{\small I primi cinque esercizi offrono una guida esplicita; i successivi cinque consolidano il calcolo; gli ultimi cinque collegano più concetti. Tutti conservano uno svolgimento completo.}\par
